\documentclass[12pt,a4paper]{article}

\usepackage[T2A]{fontenc}
\usepackage[utf8]{inputenc}
\usepackage[russian]{babel}
\usepackage{amsmath,amssymb,amsthm}
\usepackage{enumitem}
\usepackage{array}
\usepackage{listings}
\usepackage{xcolor}
\usepackage{geometry}

\geometry{margin=2.2cm}

\newtheorem{definition}{Определение}
\newtheorem{theorem}{Теорема}
\newtheorem{lemma}{Лемма}
\newtheorem{example}{Пример}
\newtheorem{remark}{Замечание}

\lstset{
  basicstyle=\ttfamily\small,
  keywordstyle=\color{blue},
  commentstyle=\color{gray},
  stringstyle=\color{red!60!black},
  breaklines=true,
  frame=single,
  columns=fullflexible
}

\title{Билет №37 \\ \small Автоматическое построение лексических и синтаксических анализаторов по формальным описаниям грамматик. Инструменты генерации лексических и синтаксических анализаторов. (СпБГУ)}
\author{Сериков Владислав Александрович}
\date{13 мая 2026}

\begin{document}

\maketitle
\tableofcontents
\newpage

\section{Общая схема трансляции и место анализаторов}

Компилятор или интерпретатор обычно содержит следующие фазы:
\[
\text{исходный текст}
\longrightarrow
\text{лексический анализ}
\longrightarrow
\text{синтаксический анализ}
\longrightarrow
\text{семантический анализ}
\longrightarrow
\text{промежуточное представление}
\longrightarrow
\text{оптимизация}
\longrightarrow
\text{кодогенерация}.
\]

Лексический анализатор, или сканер, разбивает входной поток символов на токены:
\[
\texttt{if (x == 10) y = y + 1;}
\]
преобразуется, например, в последовательность
\[
\texttt{IF},\ \texttt{LPAREN},\ \texttt{ID(x)},\ \texttt{EQ},\ \texttt{NUM(10)},\ \texttt{RPAREN},\ldots
\]

Синтаксический анализатор, или парсер, получает поток токенов и проверяет, принадлежит ли он языку, заданному грамматикой. Обычно он строит дерево разбора или абстрактное синтаксическое дерево.

\begin{definition}
\textbf{Лексема} --- конкретная подстрока входа, соответствующая некоторому токену.

\textbf{Токен} --- класс лексем, например \texttt{ID}, \texttt{NUMBER}, \texttt{IF}, \texttt{PLUS}.

\textbf{Атрибут токена} --- дополнительная информация о лексеме: имя идентификатора, значение числа, позиция в файле и т.п.
\end{definition}

Пример:

\[
\begin{array}{c|c|c}
\text{лексема} & \text{токен} & \text{атрибут} \\
\hline
\texttt{count} & \texttt{ID} & \texttt{count} \\
\texttt{123} & \texttt{NUM} & 123 \\
\texttt{+} & \texttt{PLUS} & - \\
\texttt{if} & \texttt{IF} & - \\
\end{array}
\]

\section{Формальные языки и грамматики}

\begin{definition}
\textbf{Алфавит} $\Sigma$ --- конечное непустое множество символов.

\textbf{Слово} над алфавитом $\Sigma$ --- конечная последовательность символов из $\Sigma$.

\textbf{Пустое слово} обозначается $\varepsilon$.

\textbf{Язык} над $\Sigma$ --- произвольное множество слов над $\Sigma$, то есть подмножество $\Sigma^*$.
\end{definition}

\begin{definition}
\textbf{Грамматика} --- четверка
\[
G = (N, \Sigma, P, S),
\]
где:
\begin{itemize}
    \item $N$ --- конечное множество нетерминалов;
    \item $\Sigma$ --- конечное множество терминалов, $N \cap \Sigma = \varnothing$;
    \item $P$ --- конечное множество продукций;
    \item $S \in N$ --- стартовый символ.
\end{itemize}
\end{definition}

Для синтаксического анализа языков программирования обычно используются контекстно-свободные грамматики.

\begin{definition}
Грамматика называется \textbf{контекстно-свободной}, если каждая продукция имеет вид
\[
A \to \alpha,
\]
где $A \in N$, $\alpha \in (N \cup \Sigma)^*$.
\end{definition}

Пример грамматики арифметических выражений:
\[
\begin{aligned}
E &\to E + T \mid T,\\
T &\to T * F \mid F,\\
F &\to (E) \mid id.
\end{aligned}
\]

Эта грамматика задает выражения с операциями сложения и умножения, причем умножение имеет больший приоритет, чем сложение.

\section{Лексический анализ}

\subsection{Регулярные выражения}

Лексические структуры большинства языков программирования описываются регулярными выражениями.

\begin{definition}
Регулярные выражения над алфавитом $\Sigma$ задаются индуктивно:
\begin{enumerate}
    \item $\varnothing$ --- регулярное выражение, задающее пустой язык;
    \item $\varepsilon$ --- регулярное выражение, задающее язык $\{\varepsilon\}$;
    \item если $a \in \Sigma$, то $a$ --- регулярное выражение, задающее язык $\{a\}$;
    \item если $r$ и $s$ --- регулярные выражения, то $(r|s)$, $(rs)$, $(r^*)$ --- регулярные выражения.
\end{enumerate}
\end{definition}

Операции:
\[
L(r|s) = L(r) \cup L(s),
\]
\[
L(rs) = \{xy \mid x \in L(r),\ y \in L(s)\},
\]
\[
L(r^*) = \bigcup_{k \ge 0} L(r)^k.
\]

Примеры:
\[
\texttt{digit} = [0-9],
\]
\[
\texttt{letter} = [a-zA-Z],
\]
\[
\texttt{id} = \texttt{letter}(\texttt{letter}|\texttt{digit})^*,
\]
\[
\texttt{number} = \texttt{digit}^+.
\]

\subsection{Конечные автоматы}

\begin{definition}
\textbf{Детерминированный конечный автомат} --- пятерка
\[
D = (Q, \Sigma, \delta, q_0, F),
\]
где:
\begin{itemize}
    \item $Q$ --- конечное множество состояний;
    \item $\Sigma$ --- алфавит;
    \item $\delta : Q \times \Sigma \to Q$ --- функция переходов;
    \item $q_0 \in Q$ --- начальное состояние;
    \item $F \subseteq Q$ --- множество допускающих состояний.
\end{itemize}
\end{definition}

\begin{definition}
\textbf{Недетерминированный конечный автомат с $\varepsilon$-переходами} --- пятерка
\[
N = (Q, \Sigma, \delta, q_0, F),
\]
где
\[
\delta : Q \times (\Sigma \cup \{\varepsilon\}) \to 2^Q.
\]
\end{definition}

Слово принимается автоматом, если существует путь из начального состояния в одно из допускающих состояний, помеченный этим словом, при этом $\varepsilon$-переходы не потребляют символов.

\section{Построение лексического анализатора}

Автоматическое построение лексического анализатора обычно имеет вид:
\[
\text{регулярные выражения}
\longrightarrow
\varepsilon\text{-НКА}
\longrightarrow
\text{ДКА}
\longrightarrow
\text{минимизированный ДКА}
\longrightarrow
\text{программа-сканер}.
\]

\subsection{Алгоритм Томпсона: регулярное выражение в НКА}

Алгоритм Томпсона строит по регулярному выражению эквивалентный $\varepsilon$-НКА.

\subsubsection*{Базовые случаи}

Для символа $a$:
\[
q_0 \xrightarrow{a} q_1.
\]

Для $\varepsilon$:
\[
q_0 \xrightarrow{\varepsilon} q_1.
\]

Для $\varnothing$:
\[
q_0 \quad q_1
\]
без переходов между ними.

\subsubsection*{Объединение}

Для $r|s$:
\[
\begin{array}{c}
q_0 \xrightarrow{\varepsilon} N(r) \xrightarrow{\varepsilon} q_f\\
q_0 \xrightarrow{\varepsilon} N(s) \xrightarrow{\varepsilon} q_f
\end{array}
\]

\subsubsection*{Конкатенация}

Для $rs$ финальное состояние автомата для $r$ соединяется $\varepsilon$-переходом с начальным состоянием автомата для $s$:
\[
N(r) \xrightarrow{\varepsilon} N(s).
\]

\subsubsection*{Итерация Клини}

Для $r^*$:
\[
q_0 \xrightarrow{\varepsilon} q_f,
\]
\[
q_0 \xrightarrow{\varepsilon} N(r),
\]
\[
N(r) \xrightarrow{\varepsilon} q_f,
\]
\[
N(r) \xrightarrow{\varepsilon} N(r).
\]

Более точно, если автомат для $r$ имеет начальное состояние $p$ и финальное состояние $q$, добавляются новые состояния $q_0$ и $q_f$ и переходы:
\[
q_0 \xrightarrow{\varepsilon} p,\quad
q_0 \xrightarrow{\varepsilon} q_f,\quad
q \xrightarrow{\varepsilon} p,\quad
q \xrightarrow{\varepsilon} q_f.
\]

\subsection{Пример построения НКА}

Пусть дано регулярное выражение:
\[
(a|b)^*abb.
\]

Оно задает все слова над $\{a,b\}$, заканчивающиеся на \texttt{abb}.

Построение:
\begin{enumerate}
    \item строятся автоматы для $a$ и $b$;
    \item строится автомат для $a|b$;
    \item применяется итерация Клини: $(a|b)^*$;
    \item последовательно конкатенируются автоматы для $a$, $b$, $b$.
\end{enumerate}

Полученный НКА можно затем детерминизировать.

\subsection{Теорема о регулярных выражениях и конечных автоматах}

\begin{theorem}
Для любого регулярного выражения $r$ существует конечный автомат $A$, такой что
\[
L(A) = L(r).
\]
\end{theorem}

\begin{proof}
Докажем структурной индукцией по построению регулярного выражения.

\textbf{База.}

\begin{enumerate}
    \item Для $r = \varnothing$ строится автомат с начальным и финальным состояниями без пути между ними. Он не принимает ни одного слова, поэтому его язык равен $\varnothing$.
    \item Для $r = \varepsilon$ строится автомат
    \[
    q_0 \xrightarrow{\varepsilon} q_f.
    \]
    Он принимает только пустое слово.
    \item Для $r = a$, где $a \in \Sigma$, строится автомат
    \[
    q_0 \xrightarrow{a} q_f.
    \]
    Он принимает ровно слово $a$.
\end{enumerate}

\textbf{Переход индукции.}

Пусть для регулярных выражений $r$ и $s$ уже построены автоматы $A_r$ и $A_s$, такие что
\[
L(A_r)=L(r), \qquad L(A_s)=L(s).
\]

Рассмотрим операции.

\begin{enumerate}
    \item Для $r|s$ строим новый автомат с новым начальным состоянием $q_0$ и новым финальным состоянием $q_f$. Из $q_0$ проводим $\varepsilon$-переходы в начальные состояния $A_r$ и $A_s$, а из финальных состояний $A_r$ и $A_s$ --- $\varepsilon$-переходы в $q_f$.

    Тогда слово принимается новым автоматом тогда и только тогда, когда оно принимается либо $A_r$, либо $A_s$. Следовательно,
    \[
    L(A)=L(r)\cup L(s)=L(r|s).
    \]

    \item Для $rs$ соединяем финальное состояние $A_r$ $\varepsilon$-переходом с начальным состоянием $A_s$.

    Тогда слово $w$ принимается тогда и только тогда, когда его можно представить в виде
    \[
    w = xy,
    \]
    где $x \in L(r)$, $y \in L(s)$. Следовательно,
    \[
    L(A)=L(r)L(s)=L(rs).
    \]

    \item Для $r^*$ добавляем новое начальное состояние $q_0$ и новое финальное состояние $q_f$; добавляем $\varepsilon$-переходы:
    \[
    q_0 \to q_f,\quad q_0 \to \text{start}(A_r),\quad
    \text{final}(A_r) \to q_f,\quad
    \text{final}(A_r) \to \text{start}(A_r).
    \]

    Переход $q_0 \to q_f$ принимает $\varepsilon$, а цикл через $A_r$ позволяет принять любую конечную конкатенацию слов из $L(r)$. Поэтому
    \[
    L(A)=L(r)^*=L(r^*).
    \]
\end{enumerate}

Таким образом, по любому регулярному выражению можно построить эквивалентный конечный автомат.
\end{proof}

\section{Детерминизация НКА}

\subsection{Понятие $\varepsilon$-замыкания}

\begin{definition}
$\varepsilon$-замыкание множества состояний $T$, обозначаемое $\varepsilon\text{-closure}(T)$, --- множество всех состояний, достижимых из состояний $T$ по нулю или более $\varepsilon$-переходам.
\end{definition}

\subsection{Алгоритм построения ДКА по НКА}

Пусть дан НКА:
\[
N = (Q, \Sigma, \delta, q_0, F).
\]

Строится ДКА:
\[
D = (Q_D, \Sigma, \delta_D, S_0, F_D),
\]
где состояния ДКА являются множествами состояний НКА.

\subsubsection*{Алгоритм}

\begin{enumerate}
    \item Вычислить начальное состояние ДКА:
    \[
    S_0 = \varepsilon\text{-closure}(\{q_0\}).
    \]
    \item Поместить $S_0$ в множество необработанных состояний.
    \item Пока есть необработанное состояние $S$:
    \begin{enumerate}
        \item для каждого символа $a \in \Sigma$ вычислить
        \[
        U = \varepsilon\text{-closure}(\operatorname{move}(S,a)),
        \]
        где
        \[
        \operatorname{move}(S,a)=\{q' \mid \exists q \in S:\ q' \in \delta(q,a)\};
        \]
        \item добавить переход
        \[
        S \xrightarrow{a} U;
        \]
        \item если $U$ еще не встречалось, добавить его в список необработанных состояний.
    \end{enumerate}
    \item Допускающими состояниями ДКА объявить все множества $S$, такие что
    \[
    S \cap F \ne \varnothing.
    \]
\end{enumerate}

\subsection{Минимальный пример детерминизации}

Пусть НКА распознает язык $a|b$:
\[
q_0 \xrightarrow{\varepsilon} q_1,\quad
q_0 \xrightarrow{\varepsilon} q_3,
\]
\[
q_1 \xrightarrow{a} q_2,\quad
q_3 \xrightarrow{b} q_4.
\]
Финальные состояния:
\[
F = \{q_2,q_4\}.
\]

Начальное состояние ДКА:
\[
S_0 = \varepsilon\text{-closure}(\{q_0\}) = \{q_0,q_1,q_3\}.
\]

Переходы:
\[
\operatorname{move}(S_0,a)=\{q_2\},
\]
\[
\operatorname{move}(S_0,b)=\{q_4\}.
\]

Следовательно:
\[
\{q_0,q_1,q_3\} \xrightarrow{a} \{q_2\},
\]
\[
\{q_0,q_1,q_3\} \xrightarrow{b} \{q_4\}.
\]

Состояния $\{q_2\}$ и $\{q_4\}$ являются допускающими.

\subsection{Теорема о детерминизации}

\begin{theorem}
Для любого недетерминированного конечного автомата с $\varepsilon$-переходами существует эквивалентный детерминированный конечный автомат.
\end{theorem}

\begin{proof}
Пусть дан НКА
\[
N=(Q,\Sigma,\delta,q_0,F).
\]
Построим ДКА методом подмножеств. Состояниями ДКА будут подмножества множества $Q$:
\[
Q_D \subseteq 2^Q.
\]

Начальное состояние:
\[
S_0 = \varepsilon\text{-closure}(\{q_0\}).
\]

Переход ДКА по символу $a$ из состояния $S$ определяется как:
\[
\delta_D(S,a)=\varepsilon\text{-closure}(\operatorname{move}(S,a)).
\]

Докажем, что после чтения слова $w$ ДКА находится ровно в множестве тех состояний НКА, которые достижимы из $q_0$ по слову $w$ с учетом $\varepsilon$-переходов.

Докажем это индукцией по длине слова $w$.

\textbf{База.} Если $w=\varepsilon$, то ДКА находится в состоянии
\[
S_0=\varepsilon\text{-closure}(\{q_0\}),
\]
что ровно соответствует всем состояниям НКА, достижимым из $q_0$ без чтения символов.

\textbf{Переход.} Пусть утверждение верно для слова $w$. Рассмотрим слово $wa$, где $a\in\Sigma$.

По предположению индукции после чтения $w$ ДКА находится в множестве $S$, состоящем из всех состояний НКА, достижимых после чтения $w$. После чтения символа $a$ НКА может перейти в состояния
\[
\operatorname{move}(S,a).
\]
После этого он может сделать любое число $\varepsilon$-переходов, поэтому множество достижимых состояний равно
\[
\varepsilon\text{-closure}(\operatorname{move}(S,a)).
\]
Именно так определяется переход ДКА:
\[
\delta_D(S,a)=\varepsilon\text{-closure}(\operatorname{move}(S,a)).
\]
Следовательно, утверждение верно для $wa$.

Теперь слово $w$ принимается НКА тогда и только тогда, когда среди достижимых после чтения $w$ состояний есть хотя бы одно финальное:
\[
S \cap F \ne \varnothing.
\]
Именно такие состояния ДКА объявлены финальными. Следовательно, ДКА принимает ровно тот же язык, что и НКА.
\end{proof}

\section{Минимизация ДКА}

После детерминизации часто выполняется минимизация ДКА.

\begin{definition}
Два состояния $p$ и $q$ ДКА называются \textbf{эквивалентными}, если для любого слова $w \in \Sigma^*$ автомат, стартуя из $p$, принимает $w$ тогда и только тогда, когда, стартуя из $q$, принимает $w$.
\end{definition}

Идея минимизации: объединить эквивалентные состояния.

\subsection{Алгоритм разбиения состояний}

Пусть дан ДКА:
\[
D=(Q,\Sigma,\delta,q_0,F).
\]

\subsubsection*{Алгоритм}

\begin{enumerate}
    \item Начальное разбиение:
    \[
    P_0 = \{F,\ Q\setminus F\}.
    \]
    \item Повторять уточнение разбиения:
    \begin{itemize}
        \item состояния $p$ и $q$ остаются в одном блоке, если для каждого символа $a \in \Sigma$ состояния $\delta(p,a)$ и $\delta(q,a)$ лежат в одном блоке текущего разбиения;
        \item иначе блок разделяется.
    \end{itemize}
    \item Когда разбиение перестает изменяться, каждый блок становится одним состоянием минимального ДКА.
    \item Начальный блок --- блок, содержащий $q_0$.
    \item Финальные блоки --- блоки, содержащие хотя бы одно финальное состояние исходного автомата.
\end{enumerate}

\section{Правила выбора токена}

В реальных генераторах лексических анализаторов применяется несколько соглашений.

\subsection{Правило максимального префикса}

Сканер выбирает самую длинную лексему, начинающуюся в текущей позиции.

Например, для входа:
\[
\texttt{ifx}
\]
и правил:
\[
\texttt{if} \mapsto IF,
\]
\[
[a-zA-Z][a-zA-Z0-9]^* \mapsto ID,
\]
будет выбран токен:
\[
ID(\texttt{ifx}),
\]
а не последовательность
\[
IF,\ ID(\texttt{x}).
\]

\subsection{Правило приоритета правил}

Если несколько правил совпали с одинаковой максимальной длиной, выбирается правило, записанное раньше.

Например:
\[
\begin{array}{l}
\texttt{if} \mapsto IF,\\
\lbrack a-zA-Z \rbrack \lbrack a-zA-Z0-9 \rbrack ^* \mapsto ID.
\end{array}
\]
Для входа \texttt{if} оба правила подходят, но выбирается \texttt{IF}, так как оно записано раньше.

\section{Генераторы лексических анализаторов}

\subsection{Общая структура спецификации}

Типичная спецификация для генератора лексического анализатора имеет вид:
\[
\text{определения}
\quad
%%
\quad
\text{правила}
\quad
%%
\quad
\text{пользовательский код}.
\]

Например, для \texttt{flex}:

\begin{lstlisting}
%{
#include "parser.tab.h"
%}

DIGIT   [0-9]
LETTER  [a-zA-Z_]

%%
"if"              { return IF; }
"while"           { return WHILE; }
{LETTER}({LETTER}|{DIGIT})*  { return ID; }
{DIGIT}+          { return NUM; }
"+"               { return PLUS; }
"*"               { return STAR; }
[ \t\n]+          { /* skip whitespace */ }
.                 { return ERROR; }
%%
\end{lstlisting}

Здесь каждая строка в секции правил имеет форму:
\[
\text{регулярное выражение}\quad \{\text{действие}\}.
\]

Действие выполняется, когда сканер распознает соответствующую лексему.

\subsection{Что делает генератор лексического анализатора}

Генератор лексического анализатора:
\begin{enumerate}
    \item читает регулярные выражения из спецификации;
    \item строит для каждого правила НКА;
    \item объединяет НКА всех правил в один НКА с новым начальным состоянием;
    \item детерминизирует НКА;
    \item минимизирует или оптимизирует ДКА;
    \item генерирует таблицы переходов или код с переходами;
    \item вставляет пользовательские действия в соответствующие принимающие состояния.
\end{enumerate}

\section{Синтаксический анализ}

\subsection{Дерево разбора}

\begin{definition}
\textbf{Дерево разбора} для слова $w$ по грамматике $G$ --- дерево, в котором:
\begin{itemize}
    \item корень помечен стартовым символом;
    \item внутренние вершины помечены нетерминалами;
    \item листья помечены терминалами или $\varepsilon$;
    \item если вершина помечена $A$, а ее дети слева направо помечены $X_1,\ldots,X_k$, то в грамматике есть продукция
    \[
    A \to X_1\ldots X_k.
    \]
\end{itemize}
\end{definition}

\subsection{Левый и правый вывод}

\begin{definition}
\textbf{Левый вывод} --- вывод, в котором на каждом шаге заменяется самый левый нетерминал.

\textbf{Правый вывод} --- вывод, в котором на каждом шаге заменяется самый правый нетерминал.
\end{definition}

Пример:
\[
E \Rightarrow E+T \Rightarrow T+T \Rightarrow id+T \Rightarrow id+id.
\]

\subsection{Неоднозначность грамматики}

\begin{definition}
Грамматика называется \textbf{неоднозначной}, если существует слово, имеющее два различных дерева разбора.
\end{definition}

Классический пример:
\[
E \to E+E \mid E*E \mid (E) \mid id.
\]

Для слова
\[
id+id*id
\]
существуют два дерева разбора:
\[
(id+id)*id
\]
и
\[
id+(id*id).
\]

Чтобы устранить неоднозначность, вводят уровни приоритета:
\[
\begin{aligned}
E &\to E+T \mid T,\\
T &\to T*F \mid F,\\
F &\to (E) \mid id.
\end{aligned}
\]

\section{LL-анализ}

\subsection{Идея LL-анализатора}

LL-анализатор:
\begin{itemize}
    \item читает вход слева направо;
    \item строит левый вывод;
    \item выбирает продукцию по текущему нетерминалу и одному или нескольким символам предпросмотра.
\end{itemize}

LL(1)-анализатор использует один символ предпросмотра.

\subsection{Множества FIRST и FOLLOW}

\begin{definition}
Для строки грамматических символов $\alpha$ множество $FIRST(\alpha)$ --- множество терминалов, с которых могут начинаться строки, выводимые из $\alpha$. Если из $\alpha$ выводится пустая строка, то $\varepsilon \in FIRST(\alpha)$.
\end{definition}

\begin{definition}
Для нетерминала $A$ множество $FOLLOW(A)$ --- множество терминалов, которые могут непосредственно следовать за $A$ в некоторой сентенциальной форме. Для стартового символа $S$ в $FOLLOW(S)$ добавляется маркер конца входа $\$$.
\end{definition}

\subsection{Алгоритм вычисления FIRST}

\begin{enumerate}
    \item Для терминала $a$:
    \[
    FIRST(a)=\{a\}.
    \]
    \item Для $\varepsilon$:
    \[
    FIRST(\varepsilon)=\{\varepsilon\}.
    \]
    \item Для каждого нетерминала $A$ сначала положить:
    \[
    FIRST(A)=\varnothing.
    \]
    \item Для каждой продукции
    \[
    A \to X_1X_2\ldots X_k
    \]
    добавить в $FIRST(A)$ все терминалы из $FIRST(X_1)$, кроме $\varepsilon$.
    \item Если $\varepsilon \in FIRST(X_1)$, добавить терминалы из $FIRST(X_2)$, кроме $\varepsilon$, и так далее.
    \item Если для всех $i$ выполнено $\varepsilon \in FIRST(X_i)$, добавить $\varepsilon$ в $FIRST(A)$.
    \item Повторять до неподвижной точки.
\end{enumerate}

\subsection{Алгоритм вычисления FOLLOW}

\begin{enumerate}
    \item Для стартового символа $S$ добавить:
    \[
    \$ \in FOLLOW(S).
    \]
    \item Для каждой продукции вида
    \[
    A \to \alpha B \beta
    \]
    добавить
    \[
    FIRST(\beta)\setminus\{\varepsilon\}
    \]
    в $FOLLOW(B)$.
    \item Если
    \[
    A \to \alpha B \beta
    \]
    и
    \[
    \varepsilon \in FIRST(\beta),
    \]
    то добавить
    \[
    FOLLOW(A)
    \]
    в $FOLLOW(B)$.
    \item Если продукция имеет вид
    \[
    A \to \alpha B,
    \]
    то добавить
    \[
    FOLLOW(A)
    \]
    в $FOLLOW(B)$.
    \item Повторять до неподвижной точки.
\end{enumerate}

\subsection{Построение LL(1)-таблицы}

Для каждой продукции
\[
A \to \alpha
\]
выполняются действия:
\begin{enumerate}
    \item Для каждого терминала
    \[
    a \in FIRST(\alpha)\setminus\{\varepsilon\}
    \]
    занести продукцию $A \to \alpha$ в ячейку
    \[
    M[A,a].
    \]
    \item Если
    \[
    \varepsilon \in FIRST(\alpha),
    \]
    то для каждого
    \[
    b \in FOLLOW(A)
    \]
    занести продукцию $A \to \alpha$ в ячейку
    \[
    M[A,b].
    \]
\end{enumerate}

Если в какую-либо ячейку нужно занести более одной продукции, грамматика не является LL(1).

\subsection{Критерий LL(1)}

\begin{theorem}
Грамматика является LL(1)-грамматикой тогда и только тогда, когда для любых двух различных продукций одного нетерминала
\[
A \to \alpha,\qquad A \to \beta
\]
выполнены условия:
\begin{enumerate}
    \item
    \[
    FIRST(\alpha)\cap FIRST(\beta)=\varnothing;
    \]
    \item если $\varepsilon \in FIRST(\alpha)$, то
    \[
    FIRST(\beta)\cap FOLLOW(A)=\varnothing;
    \]
    \item если $\varepsilon \in FIRST(\beta)$, то
    \[
    FIRST(\alpha)\cap FOLLOW(A)=\varnothing.
    \]
\end{enumerate}
\end{theorem}

\begin{proof}
LL(1)-анализатор выбирает продукцию по паре:
\[
(\text{верхушка стека},\ \text{текущий входной символ}).
\]
Пусть верхушка стека --- нетерминал $A$.

Если существуют две продукции
\[
A\to\alpha,\qquad A\to\beta,
\]
и
\[
FIRST(\alpha)\cap FIRST(\beta)\ne\varnothing,
\]
то существует терминал $a$, с которого может начинаться строка, выводимая и из $\alpha$, и из $\beta$. Следовательно, при предпросмотре $a$ анализатор не сможет однозначно выбрать продукцию. Значит, грамматика не LL(1).

Теперь пусть $\varepsilon \in FIRST(\alpha)$. Это означает, что $\alpha$ может исчезнуть. Тогда выбор продукции $A\to\alpha$ возможен, когда текущий входной символ принадлежит $FOLLOW(A)$, то есть следует за $A$ в некоторой выводимой строке. Если при этом существует
\[
a \in FIRST(\beta)\cap FOLLOW(A),
\]
то по символу $a$ можно выбрать и продукцию $A\to\alpha$, потому что $\alpha$ может дать $\varepsilon$, и продукцию $A\to\beta$, потому что $\beta$ может начинаться с $a$. Возникает конфликт.

Аналогично доказывается условие для случая $\varepsilon \in FIRST(\beta)$.

Обратно, если все три условия выполнены, то для каждой продукции $A\to\alpha$ множества входных символов, при которых она помещается в LL(1)-таблицу, попарно не пересекаются. Следовательно, в каждой ячейке таблицы находится не более одной продукции, и анализатор всегда делает однозначный выбор. Значит, грамматика является LL(1).
\end{proof}

\subsection{Пример LL(1)-анализа}

Грамматика:
\[
\begin{aligned}
E &\to T E',\\
E' &\to + T E' \mid \varepsilon,\\
T &\to id.
\end{aligned}
\]

Вычислим:
\[
FIRST(E)=\{id\},
\]
\[
FIRST(E')=\{+,\varepsilon\},
\]
\[
FIRST(T)=\{id\}.
\]

Множества FOLLOW:
\[
FOLLOW(E)=\{\$\},
\]
\[
FOLLOW(E')=\{\$\},
\]
\[
FOLLOW(T)=\{+,\$\}.
\]

LL(1)-таблица:
\[
\begin{array}{c|ccc}
 & id & + & \$\\
\hline
E & E\to TE' & & \\
E' & & E'\to +TE' & E'\to \varepsilon\\
T & T\to id & & \\
\end{array}
\]

Разбор строки:
\[
id+id.
\]

\[
\begin{array}{c|c|c}
\text{стек} & \text{вход} & \text{действие}\\
\hline
E\$ & id+id\$ & E\to TE'\\
TE'\$ & id+id\$ & T\to id\\
idE'\$ & id+id\$ & совпадение\ id\\
E'\$ & +id\$ & E'\to +TE'\\
+TE'\$ & +id\$ & совпадение\ +\\
TE'\$ & id\$ & T\to id\\
idE'\$ & id\$ & совпадение\ id\\
E'\$ & \$ & E'\to \varepsilon\\
\$ & \$ & принять\\
\end{array}
\]

\section{Преобразования грамматик для LL-анализа}

\subsection{Устранение левой рекурсии}

\begin{definition}
Грамматика имеет \textbf{левую рекурсию}, если существует нетерминал $A$, такой что
\[
A \Rightarrow^+ A\alpha.
\]
\end{definition}

Непосредственная левая рекурсия имеет вид:
\[
A \to A\alpha_1 \mid A\alpha_2 \mid \ldots \mid \beta_1 \mid \beta_2 \mid \ldots,
\]
где $\beta_i$ не начинаются с $A$.

Преобразование:
\[
A \to \beta_1 A' \mid \beta_2 A' \mid \ldots,
\]
\[
A' \to \alpha_1 A' \mid \alpha_2 A' \mid \ldots \mid \varepsilon.
\]

Пример:
\[
E \to E+T \mid T.
\]

После устранения:
\[
E \to T E',
\]
\[
E' \to +T E' \mid \varepsilon.
\]

\subsection{Левая факторизация}

Если грамматика содержит:
\[
A \to \alpha\beta_1 \mid \alpha\beta_2,
\]
то LL-анализатор не может выбрать продукцию после чтения начала $\alpha$.

Преобразование:
\[
A \to \alpha A',
\]
\[
A' \to \beta_1 \mid \beta_2.
\]

Пример:
\[
S \to if\ E\ then\ S\ else\ S
\mid if\ E\ then\ S.
\]

Факторизация:
\[
S \to if\ E\ then\ S\ S',
\]
\[
S' \to else\ S \mid \varepsilon.
\]

\section{LR-анализ}

\subsection{Идея LR-анализатора}

LR-анализатор:
\begin{itemize}
    \item читает вход слева направо;
    \item строит правый вывод в обратном порядке;
    \item работает снизу вверх;
    \item использует стек состояний;
    \item принимает решения на основе текущего состояния и символа предпросмотра.
\end{itemize}

LR-анализ мощнее LL-анализa и широко используется в генераторах вроде \texttt{yacc}, \texttt{bison} и аналогичных.

\subsection{Расширенная грамматика}

Для грамматики со стартовым символом $S$ вводится новый стартовый символ $S'$:
\[
S' \to S.
\]

Цель анализа --- получить конфигурацию, соответствующую завершению продукции:
\[
S' \to S \cdot
\]
при входном символе $\$$.

\subsection{LR(0)-пункты}

\begin{definition}
LR(0)-пункт --- продукция с точкой в правой части:
\[
A \to \alpha \cdot \beta.
\]
Точка показывает, какая часть продукции уже распознана.
\end{definition}

Например, для продукции:
\[
E \to E + T
\]
существуют пункты:
\[
E \to \cdot E + T,
\]
\[
E \to E \cdot + T,
\]
\[
E \to E + \cdot T,
\]
\[
E \to E + T \cdot.
\]

\subsection{Операция CLOSURE}

Если в множестве пунктов есть пункт
\[
A \to \alpha \cdot B \beta,
\]
где $B$ --- нетерминал, то нужно добавить все пункты:
\[
B \to \cdot \gamma
\]
для всех продукций $B \to \gamma$.

\subsubsection*{Алгоритм CLOSURE}

\begin{enumerate}
    \item Пусть $I$ --- исходное множество пунктов.
    \item Положить:
    \[
    CLOSURE(I)=I.
    \]
    \item Пока можно добавить новый пункт:
    \begin{itemize}
        \item если
        \[
        A\to \alpha\cdot B\beta
        \]
        уже содержится в $CLOSURE(I)$,
        \item и есть продукция
        \[
        B\to\gamma,
        \]
        \item добавить пункт
        \[
        B\to\cdot\gamma.
        \]
    \end{itemize}
    \item Вернуть полученное множество.
\end{enumerate}

\subsection{Операция GOTO}

\[
GOTO(I,X)
\]
--- множество пунктов, получаемое после распознавания символа $X$.

Формально:
\[
GOTO(I,X)=CLOSURE(\{A\to\alpha X\cdot\beta
\mid A\to\alpha\cdot X\beta \in I\}).
\]

\subsection{Построение канонического множества LR(0)-состояний}

\subsubsection*{Алгоритм}

\begin{enumerate}
    \item Построить расширенную грамматику:
    \[
    S' \to S.
    \]
    \item Начальное состояние:
    \[
    I_0 = CLOSURE(\{S'\to\cdot S\}).
    \]
    \item Для каждого уже построенного состояния $I$ и каждого символа $X \in N\cup\Sigma$ вычислить
    \[
    GOTO(I,X).
    \]
    \item Если результат непуст и ранее не встречался, добавить его как новое состояние.
    \item Повторять до неподвижной точки.
\end{enumerate}

\section{SLR(1)-таблицы}

SLR(1)-анализ использует LR(0)-состояния и множества FOLLOW.

\subsection{Построение ACTION и GOTO}

Для каждого состояния $I_i$:

\begin{enumerate}
    \item Если
    \[
    A \to \alpha \cdot a \beta \in I_i,
    \]
    где $a$ --- терминал, и
    \[
    GOTO(I_i,a)=I_j,
    \]
    то:
    \[
    ACTION[i,a]=shift\ j.
    \]

    \item Если
    \[
    A \to \alpha \cdot \in I_i,
    \]
    где $A \ne S'$, то для всех
    \[
    a \in FOLLOW(A)
    \]
    положить:
    \[
    ACTION[i,a]=reduce\ A\to\alpha.
    \]

    \item Если
    \[
    S'\to S\cdot \in I_i,
    \]
    то:
    \[
    ACTION[i,\$]=accept.
    \]

    \item Если
    \[
    GOTO(I_i,A)=I_j
    \]
    для нетерминала $A$, то:
    \[
    GOTO[i,A]=j.
    \]
\end{enumerate}

Если возникает конфликт, грамматика не является SLR(1).

\subsection{Типы конфликтов}

\begin{itemize}
    \item \textbf{Shift/reduce conflict} --- в одной ячейке нужно одновременно сделать перенос и свертку.
    \item \textbf{Reduce/reduce conflict} --- в одной ячейке нужно выполнить две разные свертки.
\end{itemize}

Пример shift/reduce-конфликта возникает в задаче висячего \texttt{else}:
\[
S \to if\ E\ then\ S
\mid if\ E\ then\ S\ else\ S
\mid other.
\]

\section{LR(1) и LALR(1)}

\subsection{LR(1)-пункты}

\begin{definition}
LR(1)-пункт имеет вид:
\[
[A \to \alpha \cdot \beta,\ a],
\]
где $a$ --- символ предпросмотра.
\end{definition}

Смысл: продукцию $A\to\alpha\beta$ можно сворачивать только тогда, когда после $A$ ожидается символ $a$.

\subsection{CLOSURE для LR(1)}

Если в множестве есть пункт:
\[
[A\to\alpha\cdot B\beta,\ a],
\]
то для каждой продукции:
\[
B\to\gamma
\]
и каждого
\[
b \in FIRST(\beta a)
\]
добавляется пункт:
\[
[B\to\cdot\gamma,\ b].
\]

\subsection{Сравнение LR-семейств}

\[
\begin{array}{c|c|c}
\text{Метод} & \text{Сила} & \text{Размер таблиц}\\
\hline
LR(0) & \text{наименьшая} & \text{малый}\\
SLR(1) & \text{больше LR(0)} & \text{малый}\\
LALR(1) & \text{почти как LR(1)} & \text{умеренный}\\
LR(1) & \text{наибольшая} & \text{большой}\\
\end{array}
\]

LALR(1) получается из LR(1) путем объединения состояний с одинаковым LR(0)-ядром. При этом таблицы обычно близки по размеру к SLR(1), но метод существенно мощнее SLR(1).

\section{Пример построения LR(0)-автомата}

Рассмотрим грамматику:
\[
S' \to S,
\]
\[
S \to a.
\]

Начальное состояние:
\[
I_0=CLOSURE(\{S'\to\cdot S\}).
\]

Так как после точки стоит нетерминал $S$, добавляем:
\[
S\to\cdot a.
\]

Итак:
\[
I_0 = \{S'\to\cdot S,\ S\to\cdot a\}.
\]

Переходы:
\[
GOTO(I_0,S)=\{S'\to S\cdot\}=I_1,
\]
\[
GOTO(I_0,a)=\{S\to a\cdot\}=I_2.
\]

Состояния:
\[
I_0=\{S'\to\cdot S,\ S\to\cdot a\},
\]
\[
I_1=\{S'\to S\cdot\},
\]
\[
I_2=\{S\to a\cdot\}.
\]

Таблица:
\[
\begin{array}{c|cc|c}
 & a & \$ & S\\
\hline
0 & s2 & & 1\\
1 & & acc & \\
2 & & r(S\to a) & \\
\end{array}
\]

Разбор входа:
\[
a\$.
\]

\[
\begin{array}{c|c|c}
\text{стек состояний} & \text{вход} & \text{действие}\\
\hline
0 & a\$ & shift\ 2\\
0\ 2 & \$ & reduce\ S\to a\\
0\ 1 & \$ & accept\\
\end{array}
\]

\section{Автоматическое построение синтаксических анализаторов}

\subsection{Общая схема}

Генератор синтаксического анализатора получает на вход:
\begin{itemize}
    \item описание токенов;
    \item контекстно-свободную грамматику;
    \item семантические действия;
    \item директивы приоритета и ассоциативности;
    \item настройки целевого языка.
\end{itemize}

На выходе он строит:
\begin{itemize}
    \item таблицы анализа или код рекурсивного спуска;
    \item обработчик ошибок;
    \item интерфейс с лексическим анализатором;
    \item иногда --- построитель дерева разбора или AST.
\end{itemize}

\subsection{Этапы работы генератора парсеров}

\begin{enumerate}
    \item Чтение грамматики.
    \item Проверка корректности спецификации.
    \item Вычисление множеств FIRST и FOLLOW.
    \item Построение автоматов состояний:
    \begin{itemize}
        \item LL-генератор строит таблицу предсказаний;
        \item LR-генератор строит каноническое множество состояний и таблицы ACTION/GOTO.
    \end{itemize}
    \item Обнаружение конфликтов.
    \item Применение директив приоритета и ассоциативности.
    \item Генерация кода.
    \item Связывание семантических действий с правилами грамматики.
\end{enumerate}

\section{Семантические действия}

Во многих генераторах грамматика дополняется фрагментами кода.

Например, в стиле \texttt{bison}:

\begin{lstlisting}
expr:
      expr '+' term   { $$ = $1 + $3; }
    | term            { $$ = $1; }
    ;

term:
      NUM             { $$ = $1; }
    ;
\end{lstlisting}

Здесь:
\begin{itemize}
    \item \texttt{\$\$} --- значение левой части продукции;
    \item \texttt{\$1}, \texttt{\$2}, \texttt{\$3} --- значения символов правой части.
\end{itemize}

Для продукции:
\[
expr \to expr + term
\]
действие:
\[
\texttt{\$\$ = \$1 + \$3}
\]
вычисляет значение выражения.

\section{Построение AST}

Дерево разбора часто содержит слишком много технических узлов. Поэтому строят абстрактное синтаксическое дерево.

Для выражения:
\[
1 + 2 * 3
\]
дерево разбора отражает все нетерминалы грамматики, а AST можно записать как:
\[
+(1,\ *(2,3)).
\]

В семантических действиях это может выглядеть так:

\begin{lstlisting}
expr:
      expr '+' term   { $$ = new Binary("+", $1, $3); }
    | term            { $$ = $1; }
    ;

term:
      term '*' factor { $$ = new Binary("*", $1, $3); }
    | factor          { $$ = $1; }
    ;

factor:
      NUM             { $$ = new Number($1); }
    ;
\end{lstlisting}

\section{Обработка ошибок}

\subsection{Лексические ошибки}

Лексическая ошибка возникает, когда ни одно регулярное выражение не подходит к текущему входу.

Пример:
\[
\texttt{x = 10 @ 20}
\]
если символ \texttt{@} не определен в языке.

Типичные стратегии:
\begin{itemize}
    \item выдать сообщение и завершить анализ;
    \item пропустить ошибочный символ;
    \item перейти к следующей строке;
    \item сформировать специальный токен ошибки.
\end{itemize}

\subsection{Синтаксические ошибки}

Синтаксическая ошибка возникает, когда текущий токен невозможен в данном состоянии парсера.

Стратегии:
\begin{itemize}
    \item аварийное завершение;
    \item panic mode: пропуск токенов до синхронизирующего множества;
    \item вставка или удаление токена;
    \item специальные правила ошибок в грамматике.
\end{itemize}

Пример правила ошибки в стиле \texttt{bison}:
\begin{lstlisting}
stmt:
      expr ';'
    | error ';'   { yyerrok; }
    ;
\end{lstlisting}

\section{Инструменты генерации анализаторов}

\subsection{Lex и Flex}

\texttt{lex} --- классический генератор лексических анализаторов.

\texttt{flex} --- современный распространенный генератор сканеров, совместимый по идее с \texttt{lex}.

Характеристики:
\begin{itemize}
    \item вход --- регулярные выражения с действиями;
    \item выход --- программа-сканер, обычно на C или C++;
    \item часто используется совместно с \texttt{yacc} или \texttt{bison};
    \item реализует правило максимального префикса и приоритет более раннего правила.
\end{itemize}

\subsection{Yacc и Bison}

\texttt{yacc} --- классический генератор синтаксических анализаторов.

\texttt{bison} --- GNU-реализация и развитие идей \texttt{yacc}.

Характеристики:
\begin{itemize}
    \item вход --- контекстно-свободная грамматика с действиями;
    \item выход --- парсер;
    \item поддерживает LALR(1), IELR(1), canonical LR(1), GLR;
    \item умеет обнаруживать конфликты shift/reduce и reduce/reduce;
    \item позволяет задавать приоритеты и ассоциативность операторов.
\end{itemize}

\subsection{ANTLR}

ANTLR --- генератор лексических и синтаксических анализаторов для структурированного текста.

Характеристики:
\begin{itemize}
    \item использует грамматики в расширенной BNF-подобной форме;
    \item строит лексер и парсер;
    \item ориентирован на LL-подобный адаптивный анализ;
    \item поддерживает генерацию кода для разных языков;
    \item широко применяется для DSL, компиляторов, анализаторов кода.
\end{itemize}

Пример грамматики ANTLR:

\begin{lstlisting}
grammar Expr;

expr
    : expr '*' expr
    | expr '+' expr
    | INT
    | '(' expr ')'
    ;

INT : [0-9]+ ;
WS  : [ \t\r\n]+ -> skip ;
\end{lstlisting}

\subsection{JavaCC}

JavaCC --- генератор синтаксических анализаторов для Java.

Особенности:
\begin{itemize}
    \item генерирует рекурсивный спуск;
    \item использует LL-подход;
    \item лексические и синтаксические правила находятся в одной спецификации;
    \item хорошо подходит для проектов, написанных на Java.
\end{itemize}

\subsection{Menhir}

Menhir --- генератор LR-парсеров для OCaml.

Особенности:
\begin{itemize}
    \item совместим по идее с \texttt{ocamlyacc};
    \item поддерживает современные средства диагностики ошибок;
    \item используется в проектах на OCaml.
\end{itemize}

\subsection{Другие инструменты}

\[
\begin{array}{c|c|c}
\text{Инструмент} & \text{Подход} & \text{Типичное применение}\\
\hline
flex & \text{регулярные выражения, ДКА} & \text{лексический анализ}\\
bison & \text{LALR/LR/GLR} & \text{синтаксический анализ}\\
ANTLR & \text{LL-подобный адаптивный анализ} & \text{лексер + парсер}\\
JavaCC & \text{LL, рекурсивный спуск} & \text{Java-парсеры}\\
Menhir & \text{LR} & \text{OCaml-парсеры}\\
Ragel & \text{конечные автоматы} & \text{протоколы, лексеры}\\
re2c & \text{ДКА} & \text{быстрые сканеры}\\
\end{array}
\]

\section{Минимальный пример flex + bison}

\subsection{Файл лексера \texttt{calc.l}}

\begin{lstlisting}
%{
#include "calc.tab.h"
%}

%%
[0-9]+      { yylval = atoi(yytext); return NUM; }
"+"         { return PLUS; }
"*"         { return STAR; }
"("         { return LPAREN; }
")"         { return RPAREN; }
[ \t\n]+    { /* skip */ }
.           { return ERROR; }
%%
\end{lstlisting}

\subsection{Файл парсера \texttt{calc.y}}

\begin{lstlisting}
%token NUM PLUS STAR LPAREN RPAREN ERROR

%%
expr:
      expr PLUS term   { $$ = $1 + $3; }
    | term             { $$ = $1; }
    ;

term:
      term STAR factor { $$ = $1 * $3; }
    | factor           { $$ = $1; }
    ;

factor:
      NUM              { $$ = $1; }
    | LPAREN expr RPAREN { $$ = $2; }
    ;
%%
\end{lstlisting}

Эта грамматика задает приоритет:
\[
* \quad \text{выше, чем} \quad +.
\]

Также она задает левую ассоциативность:
\[
1+2+3 = (1+2)+3.
\]

\section{Сравнение ручного и автоматического построения}

\[
\begin{array}{p{5cm}|p{5cm}|p{5cm}}
\text{Критерий} & \text{Ручной анализатор} & \text{Генерируемый анализатор}\\
\hline
Скорость разработки & Может быть низкой & Обычно высокая\\
Контроль над кодом & Максимальный & Ограничен шаблонами генератора\\
Диагностика ошибок & Можно сделать очень качественной & Зависит от инструмента\\
Поддерживаемость грамматики & Сложнее при росте языка & Обычно проще\\
Формальная проверяемость & Ниже & Выше, есть конфликты и отчеты\\
Производительность & Может быть очень высокой & Обычно достаточно высокая\\
\end{array}
\]

\section{Типичные проблемы при автоматической генерации}

\subsection{Неоднозначность}

Если грамматика неоднозначна, генератор может выдать конфликты.

Пример:
\[
E \to E+E \mid E*E \mid id.
\]

Решение:
\begin{itemize}
    \item переписать грамматику с уровнями приоритетов;
    \item задать приоритеты и ассоциативность в генераторе.
\end{itemize}

\subsection{Левая рекурсия}

Для LL-генераторов левая рекурсия обычно проблематична:
\[
E \to E+T \mid T.
\]

Решение:
\[
E \to T E',
\]
\[
E' \to +T E' \mid \varepsilon.
\]

Для LR-генераторов левая рекурсия, наоборот, естественна и часто предпочтительна.

\subsection{Общие префиксы}

Проблема:
\[
A \to \alpha\beta \mid \alpha\gamma.
\]

Решение:
\[
A \to \alpha A',
\]
\[
A' \to \beta \mid \gamma.
\]

\subsection{Конфликты приоритетов}

Грамматика:
\[
E \to E+E \mid E*E \mid id.
\]

Можно задать:
\[
* > +.
\]

В \texttt{bison}-подобном синтаксисе:
\begin{lstlisting}
%left PLUS
%left STAR
\end{lstlisting}

Чем позже объявлен оператор, тем выше его приоритет.

\section{Итоговая схема автоматического построения}

\subsection{Лексический анализатор}

\[
\boxed{
\text{регулярные выражения}
\to
\text{НКА}
\to
\text{ДКА}
\to
\text{таблицы переходов}
\to
\text{сканер}
}
\]

\subsection{Синтаксический анализатор}

\[
\boxed{
\text{КС-грамматика}
\to
\text{FIRST/FOLLOW или LR-состояния}
\to
\text{таблицы анализа}
\to
\text{парсер}
}
\]

\subsection{Основные идеи}

\begin{enumerate}
    \item Лексика языков программирования обычно описывается регулярными выражениями.
    \item Регулярные выражения автоматически преобразуются в конечные автоматы.
    \item Синтаксис обычно описывается контекстно-свободными грамматиками.
    \item LL-анализ строит левый вывод сверху вниз.
    \item LR-анализ строит правый вывод в обратном порядке снизу вверх.
    \item Генераторы анализаторов автоматизируют построение таблиц, обнаружение конфликтов и генерацию кода.
    \item Качество грамматики существенно влияет на качество полученного анализатора.
\end{enumerate}

\end{document}
